\section{Relational ensembles}

\begin{frame}{The relational partition function}
  \bookref{Chapter 2, Section 1}
  Let $E(S)$ be isomorphism-invariant. Summing over isomorphism classes,
  \[
    \boxed{
      Z(\beta)=\sum_{[S]}
      \frac{n!}{|\Aut(S)|}e^{-\beta E(S)}
      =\sum_{[S]}e^{I(S)-\beta E(S)}.}
  \]

  The induced Gibbs law is
  \[
    \mathbb P_\beta([S])
      =\frac{e^{I(S)-\beta E(S)}}{Z(\beta)}.
  \]

  \begin{block}{Interpretation}
    Energy weights structures; each orbit counts labeled realizations.
    $X$ labels entities; configurations are structures on $X$.
  \end{block}
  Physical use requires distinct admissible configurations. Identify redundant
  labels first and specify the background-preserving reference action.
\end{frame}

\begin{frame}{Energy--multiplicity competition}
  \bookref{Chapter 2, Sections 1, 2, and 5}
  \begin{columns}[T]
    \column{0.48\textwidth}
      \begin{block}{Highly symmetric type}
        \begin{itemize}
          \item large $|\Aut(S)|$,
          \item small $I(S)$,
          \item few labeled realizations.
        \end{itemize}
      \end{block}
    \column{0.48\textwidth}
      \begin{block}{Rigid type}
        \begin{itemize}
          \item $|\Aut(S)|=1$,
          \item $I(S)=\log(n!)$,
          \item a full orbit of $n!$ realizations.
        \end{itemize}
      \end{block}
  \end{columns}

  \vspace{0.5em}
  For singleton relations on three points, a rewarded loop contributes
  $3e^{\beta\varepsilon}$ while a directed edge contributes $6$. Thus the loop
  wins exactly when $\beta\varepsilon>\log 2$.

  \begin{alertblock}{Not thermodynamic entropy by definition}
    $I(S)$ is an orbit-counting entropy-like term. Its physical meaning depends
    on the ensemble and on which labels represent genuine microstates.
  \end{alertblock}
\end{frame}

\begin{frame}{Why asymmetry is statistically relevant}
  \bookref{Chapter 2, Section 2}
  In the classical uniform ensemble of undirected graphs, Erd\H{o}s and
  R\'enyi proved that almost all large graphs are asymmetric
  \cite{erdosrenyi1963}.

  \vspace{0.5em}
  \begin{center}
    \begin{tikzpicture}[>=Latex]
      \draw[->,very thick] (0,0)--(9,0);
      \node[below] at (0,0) {large automorphism group};
      \node[below] at (9,0) {rigid};
      \foreach \x in {0.6,1.7,3.0,4.4,5.9,7.1,7.8,8.3,8.65,8.85}
        \fill[RelBlue] (\x,0.13) circle (2.2pt);
      \node[above,align=center] at (7.5,0.45)
        {generic mass of many\\classical graph ensembles};
    \end{tikzpicture}
  \end{center}

  \begin{block}{Scope of the precedent}
    It supports the importance of rigid orbit types, but it is not a theorem
    about every relational ensemble, every energy, or law-constrained algebraic
    structures.
  \end{block}
\end{frame}
